\documentclass[11pt,a4paper]{article}

% ── Packages ──────────────────────────────────────────────────────
\usepackage[utf8]{inputenc}
\usepackage[T1]{fontenc}
\usepackage{amsmath,amssymb}
\usepackage{graphicx}
\usepackage{booktabs}
\usepackage{hyperref}
\usepackage[margin=1in]{geometry}
\usepackage{enumitem}
\usepackage{algorithm}
\usepackage{algorithmic}
\usepackage{subcaption}
\usepackage{xcolor}

\hypersetup{colorlinks=true, linkcolor=blue, citecolor=blue, urlcolor=blue}

\title{Adaptive Vision Pipeline:\\
       Classification, Detection \& Segmentation}
\author{Khaled Koubaa}
\date{\today}

\begin{document}
\maketitle

% ══════════════════════════════════════════════════════════════════
\begin{abstract}
We present an end-to-end computer-vision pipeline that unifies image
classification, object detection, and instance segmentation within a single
config-driven framework.  Built on PyTorch and OpenCV, the system leverages a
shared ResNet-50 backbone with task-specific heads—Faster~R-CNN for detection
and Mask~R-CNN for segmentation—allowing pixel-level masks to be toggled on
whenever they improve downstream accuracy.  We detail the architecture,
training procedure, evaluation protocol, and experimental results on a custom
dataset.
\end{abstract}

% ══════════════════════════════════════════════════════════════════
\section{Introduction}
\label{sec:intro}

Modern vision systems increasingly require \emph{multi-granularity}
understanding of scenes: recognising \emph{what} is in an image
(classification), \emph{where} objects are located (detection), and
\emph{which pixels} belong to each instance (segmentation).  Rather than
maintaining three independent code-bases, this project exposes a unified
pipeline whose behaviour is controlled by a single YAML configuration file.

\paragraph{Contributions.}
\begin{enumerate}[nosep]
  \item A modular Python pipeline supporting three vision tasks with a shared
        backbone.
  \item OpenCV-based preprocessing with geometric and photometric
        augmentations that respect bounding-box and mask annotations.
  \item Comprehensive evaluation: confusion matrices (classification),
        mAP/IoU curves (detection), and mask-mAP (segmentation).
  \item Reproducible experiments with deterministic seeding, mixed-precision
        training, and early stopping.
\end{enumerate}

% ══════════════════════════════════════════════════════════════════
\section{Methodology}
\label{sec:method}

\subsection{Backbone Architecture}

All three tasks share a ResNet-50~\cite{he2016deep} feature extractor
pretrained on ImageNet~\cite{deng2009imagenet}.  For detection and
segmentation, a Feature Pyramid Network (FPN)~\cite{lin2017feature} is added
to produce multi-scale feature maps.

\subsection{Classification Head}

The global average-pooled features $\mathbf{z}\in\mathbb{R}^{2048}$ are
passed through a linear classifier:
\begin{equation}
  \hat{y} = \arg\max_{c}\;\mathbf{W}_c^\top \mathbf{z} + b_c\,,
  \quad c \in \{1,\dots,C\}
  \label{eq:cls}
\end{equation}
trained with the cross-entropy loss:
\begin{equation}
  \mathcal{L}_{\text{cls}} = -\frac{1}{N}\sum_{i=1}^{N}
    \log\frac{\exp(f_{y_i}(\mathbf{x}_i))}
             {\sum_{c=1}^{C}\exp(f_c(\mathbf{x}_i))}
  \label{eq:ce}
\end{equation}

\subsection{Detection Head (Faster R-CNN)}

Faster R-CNN~\cite{ren2015faster} adds a Region Proposal Network (RPN) on
top of the FPN features.  The RPN proposes candidate bounding boxes that are
refined and classified by the RoI (Region of Interest) head.

\paragraph{RPN loss.}  Each anchor is assigned a binary label
$y^*\!\in\{0,1\}$ and a regression target
$\mathbf{t}^*\!=\!(t_x^*,t_y^*,t_w^*,t_h^*)$:
\begin{equation}
  \mathcal{L}_{\text{RPN}} =
    \frac{1}{N_{\text{cls}}}\sum_i \mathcal{L}_{\text{cls}}(p_i, y_i^*)
    + \lambda\,\frac{1}{N_{\text{reg}}}\sum_i y_i^*\,
      \text{smooth}_{L_1}(\mathbf{t}_i - \mathbf{t}_i^*)
  \label{eq:rpn}
\end{equation}

\paragraph{RoI head loss.}  Proposals are pooled to fixed-size features via
RoI Align~\cite{he2017mask} and fed to fully-connected layers for
classification and bounding-box regression.

\subsection{Segmentation Head (Mask R-CNN)}

Mask R-CNN~\cite{he2017mask} extends Faster R-CNN with a parallel mask
branch that predicts a binary mask for each detected instance:
\begin{equation}
  \mathcal{L}_{\text{mask}} =
    -\frac{1}{m^2}\sum_{1\le i,j\le m}
    \bigl[y_{ij}\log\hat{m}_{ij}
          + (1-y_{ij})\log(1-\hat{m}_{ij})\bigr]
  \label{eq:mask}
\end{equation}
where $m\times m$ is the mask resolution and $\hat{m}_{ij}$ is the predicted
probability for pixel $(i,j)$.  This head is \emph{only activated} when the
configuration flag \texttt{enable\_masks} is set to \texttt{true}.

\subsection{Evaluation Metrics}

\paragraph{Classification.}  Accuracy, macro-averaged precision, recall,
$F_1$-score, and the full confusion matrix.

\paragraph{Detection.}  Mean Average Precision at IoU threshold~$\theta$:
\begin{equation}
  \text{AP}_c = \int_0^1 p_c(r)\,dr\,,
  \qquad
  \text{mAP} = \frac{1}{C}\sum_{c=1}^{C}\text{AP}_c
  \label{eq:map}
\end{equation}
where $p_c(r)$ is the precision--recall curve for class~$c$.

\paragraph{Segmentation.}  Mask-mAP is computed identically but using mask
IoU instead of box IoU:
\begin{equation}
  \text{IoU}_{\text{mask}} =
    \frac{|M_{\text{pred}} \cap M_{\text{gt}}|}
         {|M_{\text{pred}} \cup M_{\text{gt}}|}
  \label{eq:mask_iou}
\end{equation}

% ══════════════════════════════════════════════════════════════════
\section{Experimental Setup}
\label{sec:setup}

\subsection{Hyper-parameters}

\begin{table}[h]
\centering
\caption{Training hyper-parameters per task.}
\label{tab:hyper}
\begin{tabular}{lccc}
\toprule
Parameter          & Classification & Detection & Segmentation \\
\midrule
Backbone           & ResNet-50      & ResNet-50-FPN & ResNet-50-FPN \\
Input size         & $224\times224$ & $640\times640$ & $640\times640$ \\
Batch size         & 32             & 4              & 2 \\
Optimiser          & SGD            & SGD            & SGD \\
Learning rate      & $10^{-3}$      & $5\times10^{-3}$ & $5\times10^{-3}$ \\
Weight decay       & $10^{-4}$      & $5\times10^{-4}$ & $5\times10^{-4}$ \\
LR schedule        & Cosine         & StepLR (15)    & StepLR (15) \\
Epochs             & 30             & 50             & 50 \\
Early stopping     & 10 epochs      & 15 epochs      & 15 epochs \\
Mixed precision    & Yes            & Yes            & Yes \\
\bottomrule
\end{tabular}
\end{table}

\subsection{Data Augmentation}

All augmentations are implemented in OpenCV:
\begin{itemize}[nosep]
  \item Resize to target resolution (bilinear interpolation).
  \item Random horizontal flip ($p=0.5$) with box/mask coordinate adjustment.
  \item Random brightness and contrast jitter ($\pm20\%$).
  \item ImageNet channel-wise normalisation ($\mu$, $\sigma$).
\end{itemize}

% ══════════════════════════════════════════════════════════════════
\section{Results}
\label{sec:results}

\textit{(Fill in after running experiments on target dataset.)}

\begin{figure}[h]
  \centering
  % \includegraphics[width=0.9\textwidth]{figures/training_curves.png}
  \caption{Training and validation loss/metric curves over epochs.}
  \label{fig:curves}
\end{figure}

\begin{figure}[h]
  \centering
  % \includegraphics[width=0.6\textwidth]{figures/confusion_matrix.png}
  \caption{Confusion matrix on the validation set (classification task).}
  \label{fig:cm}
\end{figure}

\begin{figure}[h]
  \centering
  % \includegraphics[width=0.9\textwidth]{figures/prediction_grid.png}
  \caption{Sample detection/segmentation predictions on validation images.}
  \label{fig:preds}
\end{figure}

% ══════════════════════════════════════════════════════════════════
\section{Discussion}
\label{sec:discuss}

The modular design allows rapid experimentation: switching from detection to
segmentation requires changing a single flag (\texttt{enable\_masks:~true})
in the configuration file.  The Mask~R-CNN head adds minimal overhead to
Faster~R-CNN while providing pixel-level instance boundaries that can
improve downstream applications such as measurement, counting, or
quality inspection.

% ══════════════════════════════════════════════════════════════════
\section{Conclusion}
\label{sec:conclusion}

We demonstrated a reproducible, config-driven pipeline for multi-task visual
understanding.  The shared ResNet-50/FPN backbone keeps inference efficient,
while task-specific heads allow flexible deployment depending on application
requirements.  All code, trained weights, and this report are provided for
full reproducibility.

% ══════════════════════════════════════════════════════════════════
\bibliographystyle{plain}
\bibliography{references}

\end{document}
